
Calculons maintenant le second membre de la formule de Riemann--Roch. Pour tout \(k\), on a :
\[
 \widetilde c(\xi^k)=(\widetilde c(\xi))^k=(1,1+x)^k=(1,1+kx)
\]
d'après (V (6.2.1)). Par suite :
\begin{equation}\tag{5.4.2}
 \ch(\xi^k)=\exp(kx).
\end{equation}
Le second membre de la formule de Riemann--Roch est donc :
\begin{equation}\tag{5.4.3}
 f_{\mathrm{gr}}\left(\exp(kx)\cdot
 \frac{x^{r+1}}{(1-\exp(-x))^{r+1}}\right)
\end{equation}
Comme \(x\) est la classe dans \(\Gr^1K^{\bullet}(P)\) de \(1-\xi^{-1}\), \(f_{\mathrm{gr}}(x^n)\) est, pour \(0\leq n\leq r\), la classe dans
\(\Gr^{n-r}K^{\bullet}(Y)\) de \(f_*((1-\xi^{-1})^n)=1\). On a donc :
\begin{equation}\tag{5.4.4}
 f_{\mathrm{gr}}(x^n)=0\quad\text{si }0\leq n<r~;\qquad
 f_{\mathrm{gr}}(x^r)=1.
\end{equation}
Rappelons enfin que les éléments \(1,x,\ldots,x^{r+1}\) sont linéairement dépendants sur \(\Gr K^{\bullet}(Y)\) (VI 5.6), et vérifient l'équation :
\begin{equation}\tag{5.4.5}
 x^{r+1}=0,
\end{equation}
et que \(1,\ldots,x^r\) sont linéairement indépendants sur \(\Gr K^{\bullet}(Y)\).

On peut alors écrire la série \(\varphi_k=\exp(kx)\cdot\frac{x^{r+1}}{(1-\exp(-x))^{r+1}}\) de façon unique comme combinaison linéaire à coefficients dans
\(\widehat{\Gr K^{\bullet}(Y)_{\mathbb Q}}\) des éléments \(1,\ldots,x^r\). L'image par \(f_{\mathrm{gr}}\) de \(\varphi_k\) est alors, d'après (5.4.4) et 3.7, son coefficient sur \(x^r\). Compte tenu de (5.4.1), le théorème de Riemann--Roch pour \(f\) est donc équivalent à :
\[
\left\{
\begin{array}{ll}
\text{pour }-r\leq k<0,\text{ le coefficient de }x^r\text{ dans }\varphi_k\text{ est }0~;\\
\text{pour }k=0,\text{ le coefficient de }x^r\text{ dans }\varphi_k\text{ est }1.
\end{array}
\right.
\tag{RR}
\]
Ce coefficient est le résidu de la forme différentielle
\[
 \frac{\exp(kx)\cdot dx}{(1-\exp(-x))^{r+1}}.
\]
